% !TeX program = pdflatex
\documentclass[aspectratio=169,xcolor={dvipsnames,table}]{beamer}

\input{../../_en_preambulo_beamer}
% Comandos y macros estadísticas compartidas para las presentaciones Beamer en inglés (Metropolis)

% Conjuntos numéricos básicos
\newcommand{\R}{\mathbb{R}}
\newcommand{\N}{\mathbb{N}}
\newcommand{\Q}{\mathbb{Q}}
\newcommand{\Z}{\mathbb{Z}}
\newcommand{\set}[1]{\left\{ #1 \right\}}
\newcommand{\abs}[1]{\left|#1\right|}
\newcommand{\norm}[1]{\left\| #1 \right\|}

% Comandos estadísticos y probabilísticos del curso
\newcommand{\Var}{\operatorname{Var}}
\newcommand{\cov}{\operatorname{Cov}}
\newcommand{\corr}{\rho}
\newcommand{\comb}[2]{\begin{pmatrix} #1 \\ #2 \end{pmatrix}}
\newcommand{\s}{\sigma}
\newcommand{\particion}{\mathfrak{P}}
\newcommand{\card}[1]{\#\left( #1 \right)}
\newcommand{\seq}[3]{\set{#1}_{#2}^{#3}}
\newcommand{\E}{\mathbb{E}}
\newcommand{\Prob}{\mathbb{P}}

% Entornos pedagógicos y cajas en inglés académico natural (soportando tanto nombres en inglés como en español)
\newenvironment{discussion}{%
  \begin{alertblock}{Discussion Question}%
}{%
  \end{alertblock}%
}
\newenvironment{discusion}{%
  \begin{alertblock}{Discussion Question}%
}{%
  \end{alertblock}%
}

\newenvironment{reflection}{%
  \begin{alertblock}{Classroom Reflection}%
}{%
  \end{alertblock}%
}
\newenvironment{reflexion}{%
  \begin{alertblock}{Classroom Reflection}%
}{%
  \end{alertblock}%
}

\newenvironment{paradox}{%
  \begin{alertblock}{Exploratory Paradox}%
}{%
  \end{alertblock}%
}
\newenvironment{paradoja}{%
  \begin{alertblock}{Exploratory Paradox}%
}{%
  \end{alertblock}%
}

\newenvironment{motivation}{%
  \begin{exampleblock}{Motivation: Data Science in Practice}%
}{%
  \end{exampleblock}%
}
\newenvironment{motivacion}{%
  \begin{exampleblock}{Motivation: Data Science in Practice}%
}{%
  \end{exampleblock}%
}

\newenvironment{quickchallenge}{%
  \begin{block}{Quick Team Challenge}%
}{%
  \end{block}%
}
\newenvironment{retoclase}{%
  \begin{block}{Quick Team Challenge}%
}{%
  \end{block}%
}


\title{Gamma, Beta, and Weibull Distributions}
\subtitle{Section 04.07 --- Waiting Times, Proportions, and Reliability}
\author[J. Castillo Colmenares]{Juliho Castillo Colmenares}
\institute[Tec de Monterrey]{Tecnologico de Monterrey}
\date{\vspace{-1.5cm}}

\begin{document}

% SLIDE 01: Title Page
\begin{frame}[plain]
	\titlepage
\end{frame}

% SLIDE 02: Roadmap
\begin{frame}{Roadmap --- Chapter 04: Continuous Random Variables}
	\vspace{-0.15cm}
	\begin{block}{Curricular Structure of Unit 3}
		\footnotesize
		\begin{enumerate}\itemsep=0.03cm
			\item \textcolor{gray}{Section 04.01: PDF and Continuous Support}
			\item \textcolor{gray}{Section 04.02: Continuous CDF and Quantiles}
			\item \textcolor{gray}{Section 04.03: Expectation, Variance, and Continuous LOTUS}
			\item \textcolor{gray}{Section 04.04: Continuous Uniform Distribution}
			\item \textcolor{gray}{Section 04.05: Normal Distribution}
			\item \textcolor{gray}{Section 04.06: Exponential Distribution and Memoryless Processes}
			\item \textbf{\textcolor{TecRojo}{Section 04.07: Gamma, Beta, and Weibull Distributions}} $\leftarrow$ \emph{Current focus --- Chapter closes}
		\end{enumerate}
	\end{block}
	\vspace{-0.2cm}
	\begin{alertblock}{Learning Objective}
		\scriptsize
		Generalize the Gamma family beyond the Exponential and Chi-squared, introduce the Beta distribution for modeling proportions on $[0,1]$, and master the Weibull distribution as the central tool for reliability analysis and time-varying failure rates.
	\end{alertblock}
\end{frame}

% SLIDE 03: Motivation
\begin{frame}{Three Flexible Shapes Beyond the Symmetric}
	\footnotesize
	\begin{motivacion}
		So far we have seen distributions with fixed shapes: uniform (flat), exponential (decreasing), and normal (symmetric). Many real phenomena require more flexible shapes: waiting times until the $\alpha$-th event, proportions bounded on $[0,1]$, and failure times whose risk changes with the component's age.
	\end{motivacion}
	\vspace{0.15cm}
	\pause
	\begin{itemize}\itemsep=0.1cm
		\item \textbf{Gamma:} Generalizes the Exponential by summing $\alpha$ waiting times; contains the Chi-squared as a special case. \pause
		\item \textbf{Beta:} The canonical distribution on $[0,1]$; natural conjugate prior for Bayesian proportions. \pause
		\item \textbf{Weibull:} Generalizes the Exponential by allowing increasing or decreasing failure rates; a pillar of reliability engineering.
	\end{itemize}
\end{frame}

% SLIDE 03B: Gamma Distribution
\begin{frame}{Gamma Distribution: Definition and Erlang Property}
	\vspace{-0.15cm}
	\begin{columns}[T]
		\begin{column}{0.48\textwidth}
			\begin{block}{PDF and Properties}
				\footnotesize
				$f(x) = \frac{1}{\Gamma(\alpha)\beta^{\alpha}}x^{\alpha-1}e^{-x/\beta}$, $x>0$. Mean $\mu_{X}=\alpha\beta$, variance $\sigma_{X}^{2}=\alpha\beta^{2}$. MGF: $M_{X}(t)=(1-\beta t)^{-\alpha}$.
			\end{block}
		\end{column}
		\begin{column}{0.48\textwidth}
			\begin{alertblock}{Special Cases and Additivity}
				\footnotesize
				$\text{Gamma}(1,\beta)=\text{Exp}(1/\beta)$; $\text{Gamma}(\nu/2,2)=\chi^{2}_{\nu}$. If $X_{i}\sim\text{Gamma}(\alpha_{i},\beta)$ independent, $\sum X_{i} \sim \text{Gamma}(\sum\alpha_{i},\beta)$ (Erlang).
			\end{alertblock}
		\end{column}
	\end{columns}
\end{frame}

% SLIDE 04: Exponential and Chi-Squared Cases
\begin{frame}{Special Cases of the Gamma Distribution}
\scriptsize
If $X\sim\Gamma(\alpha,\theta)$, then two important cases are
\[
\alpha=1\Rightarrow X\sim\operatorname{Exp}(1/\theta),
\qquad
\alpha=\nu/2,\ \theta=2\Rightarrow X\sim\chi^2_{\nu}.
\]
The exponential case models memoryless waiting times; the chi-squared case is
central to variance inference and goodness-of-fit procedures.
\end{frame}

% SLIDE 05: Beta Distribution
\begin{frame}{Beta Distribution: Definition and Bayesian Prior}
	\vspace{-0.15cm}
	\begin{columns}[T]
		\begin{column}{0.48\textwidth}
			\begin{block}{PDF and Moments}
				\footnotesize
				$f(x) = \frac{1}{B(\alpha,\beta)}x^{\alpha-1}(1-x)^{\beta-1}$, $0<x<1$, with $B(\alpha,\beta)=\Gamma(\alpha)\Gamma(\beta)/\Gamma(\alpha+\beta)$. Mean $\alpha/(\alpha+\beta)$, symmetry $\text{Beta}(\alpha,\beta)=1-\text{Beta}(\beta,\alpha)$.
			\end{block}
		\end{column}
		\begin{column}{0.48\textwidth}
			\begin{alertblock}{Conjugate Prior}
				\footnotesize
				If prior on $p$ is $\text{Beta}(\alpha,\beta)$ and $k$ successes in $n$ Bernoulli trials are observed, posterior is $\text{Beta}(\alpha+k,\beta+n-k)$. Standard for proportions and A/B testing.
			\end{alertblock}
		\end{column}
	\end{columns}
\end{frame}

% SLIDE 05B: Formal Beta Density
\begin{frame}{Formal Definition of the Beta Distribution}
\scriptsize
For $0<x<1$ and $\alpha,\beta>0$,
\[
f_X(x)=\frac{x^{\alpha-1}(1-x)^{\beta-1}}{B(\alpha,\beta)},
\qquad
B(\alpha,\beta)=\int_0^1t^{\alpha-1}(1-t)^{\beta-1}\,dt.
\]
Its support is bounded, making the Beta family natural for proportions and
probabilities.
\end{frame}

% SLIDE 06: Weibull Distribution
\begin{frame}{Weibull Distribution: Definition and Hazard Function}
	\vspace{-0.15cm}
	\begin{columns}[T]
		\begin{column}{0.48\textwidth}
			\begin{block}{PDF and CDF}
				\footnotesize
				$f(t) = \frac{\beta}{\eta}(t/\eta)^{\beta-1}e^{-(t/\eta)^{\beta}}$, $t\ge0$; $F(t)=1-e^{-(t/\eta)^{\beta}}$. Shape $\beta$, scale $\eta$.
			\end{block}
		\end{column}
		\begin{column}{0.48\textwidth}
			\begin{alertblock}{Hazard Rate $h(t)=(\beta/\eta)(t/\eta)^{\beta-1}$}
				\footnotesize
				$\beta=1$: constant rate, $T\sim\text{Exp}(1/\eta)$ (no aging). $\beta<1$: decreasing (infant mortality). $\beta>1$: increasing (wear-out).
			\end{alertblock}
		\end{column}
	\end{columns}
\end{frame}

% SLIDE 06: Applications
\begin{frame}{Applications in Engineering, Reliability, and Data Science}
	\vspace{-0.1cm}
	\begin{itemize}\itemsep=0.1cm
		\item \textbf{Industrial Reliability:} Weibull models failure times in manufacturing, predictive maintenance, and product warranties. \pause
		\item \textbf{Queueing Theory:} Gamma/Erlang models multi-stage service times ($M/E_{k}/1$). \pause
		\item \textbf{Bayesian Statistics:} Beta is the standard conjugate prior for conversion rates and error rates in A/B testing. \pause
		\item \textbf{Data Science:} Bayesian regularization, hierarchical models, and survival analysis in health and engineering.
	\end{itemize}
\end{frame}

% SLIDE 07: Python Lab Block 1
\begin{frame}[fragile]{Python Lab: Gamma Distribution and Erlang Property}
	\vspace{-0.05cm}
	\scriptsize
	Validation of $\int f = 1$ for the Gamma, the Erlang additive property, and the Exponential/Chi-squared special cases (\texttt{04.07\_gamma\_beta\_weibull.py}):
	{\lstset{basicstyle=\fontsize{3.6pt}{4.3pt}\selectfont\ttfamily,aboveskip=0.1em,belowskip=0.1em}
	\lstinputlisting[firstline=17, lastline=51]{../../code/04_variables_aleatorias_continuas/04.07_gamma_beta_weibull.py}}
\end{frame}

% SLIDE 08: Python Lab Block 2
\begin{frame}[fragile]{Python Lab: Beta Distribution and Bayesian Update}
	\vspace{-0.05cm}
	\scriptsize
	Moment calculation, symmetry property verification, and a conjugate-prior update:
	{\lstset{basicstyle=\fontsize{4pt}{5pt}\selectfont\ttfamily}
	\lstinputlisting[firstline=54, lastline=83]{../../code/04_variables_aleatorias_continuas/04.07_gamma_beta_weibull.py}}
\end{frame}

% SLIDE 09: Python Lab Block 3
\begin{frame}[fragile]{Python Lab: Weibull and Reliability Analysis}
	\vspace{-0.05cm}
	\scriptsize
	Hazard function comparison and reliability curves across shape parameters:
	{\lstset{basicstyle=\fontsize{4pt}{5pt}\selectfont\ttfamily}
	\lstinputlisting[firstline=86, lastline=115]{../../code/04_variables_aleatorias_continuas/04.07_gamma_beta_weibull.py}}
\end{frame}

% SLIDE 10: Python Lab Terminal Output
\begin{frame}[fragile]{Python Lab: Terminal Output}
	\vspace{-0.1cm}
	\begin{block}{}
		\fontsize{4pt}{5pt}\selectfont
		\begin{verbatim}
=== Block 1: Gamma Distribution & Erlang ===
Gamma(alpha=3, beta=2) PDF integral: 1.00000000
  Mean: 6.0000, Variance: 12.0000
Erlang: Gamma(3,0.3333); P(T>1.5)=0.173578
Exponential case: Gamma(1,4) vs Exp(0.25): P(X>6)=0.223130 (match)
Chi-squared case: nu=3, P(X>7.815)=0.049994 (match)

=== Block 2: Beta Distribution & Moments ===
Beta(alpha=3, beta=5) PDF integral: 1.00000000
  Mean: 0.375000, Variance: 0.026042
Symmetry: P(X<=0.3)|Beta(2,5) = 1-P(Y<=0.7)|Beta(5,2) = 0.579825
Bayesian update: prior Beta(2,2) + 14/20 -> Beta(16,8), mean=0.666667

=== Block 3: Weibull & Reliability Analysis ===
Weibull(shape=1.5, scale=4) PDF integral: 1.00000000
  R(3) = 0.522297
Hazard: Lot A (beta=1) h=0.100 constant; Lot B (beta=2) h(5)=0.100, h(15)=0.300
		\end{verbatim}
	\end{block}
\end{frame}









% SLIDE 19: Closing and Chapter Perspective
\begin{frame}{Closing Section and Chapter Perspective}
	\vspace{-0.2cm}
	\begin{block}{Synthesis: Three Flexible Continuous Families}
		\scriptsize
		The Gamma, Beta, and Weibull distributions complete our study of the fundamental continuous random variables. Gamma generalizes waiting times, Beta models proportions on $[0,1]$, and Weibull captures time-varying failure rates. Together they cover the vast majority of waiting, proportion, and reliability phenomena in practice.
	\end{block}
	\vspace{0.05cm}
	\pause
	\begin{alertblock}{Key Lessons and Next Chapter}
		\begin{itemize}\itemsep=0.05cm
			\item \textbf{Gamma:} $f(x) \propto x^{\alpha-1}e^{-x/\beta}$; generalizes Exponential and Chi-squared; additive under sums.
			\item \textbf{Beta:} $f(x) \propto x^{\alpha-1}(1-x)^{\beta-1}$ on $[0,1]$; standard Bayesian conjugate prior.
			\item \textbf{Weibull:} hazard $h(t) = (\beta/\eta)(t/\eta)^{\beta-1}$; $\beta=1$ no aging, $\beta>1$ wear-out.
			\item \textbf{Chapter 04 complete.} \textbf{Next:} \textcolor{TecRojo}{Chapter 05 --- Sampling Distributions}.
		\end{itemize}
	\end{alertblock}
\end{frame}

% SLIDE 20: Modular Perspective
\begin{frame}{Modular Perspective --- Chapter 04 Closure}
\scriptsize
\begin{alertblock}{Continuous Families Completed}
Gamma, Beta, and Weibull extend the symmetric Normal and bounded Uniform models
to waiting times, proportions, and reliability data.
\end{alertblock}
\begin{block}{Transition}
The next chapter develops sampling distributions and the inferential behavior of
statistics computed from random samples.
\end{block}
\end{frame}

\end{document}
